\documentclass[11pt,a4paper]{coverletter}

\senderblock
  {Radheem Bin Razi}
  {Distributed Systems \& Cloud-Native Engineer}
  {Ilmenau, Germany \textbar\ \href{mailto:sheikh.radheem@gmail.com}{sheikh.radheem@gmail.com} \textbar\ \href{https://radheem.github.io/my-cv}{portfolio}}

\begin{document}

%% ===== ENGLISH =====
\recipient{T-Systems International}{Hiring Team}{}

\opening{Dear Hiring Team,}

Building a cloud platform that quietly absorbs complexity so customers can focus on their actual business is exactly the kind of infrastructure work I want to do next. I am applying for the Backend Engineer T Cloud Public – Container role at T-Systems International because your focus on reducing digital transformation friction aligns with how I approach distributed systems. My background centers on designing containerized backends and automating their deployment lifecycles.

In my recent work, I have consistently replaced fragile manual processes with automated, container-native pipelines. I designed a Kubernetes deployment strategy using kustomize and Skaffold that let developers iterate on microservices without managing underlying cluster state. I also rebuilt a messaging layer by migrating away from a heavy sidecar architecture to a lighter NATS JetStream setup, which simplified cross-service communication and reduced operational overhead. These experiences directly support your need for robust CI/CD practices and infrastructure automation within a public cloud environment.

Reliability and observability have been central to every system I ship. I built a metrics pipeline that routes real-time performance data through Kafka into Grafana dashboards, giving operators clear visibility into service health before issues impact users. I also implemented strict schema validation and automated quality gates that catch data drift and configuration errors early in the release cycle. This focus on proactive monitoring and consistent testing matches your emphasis on developing monitoring and quality assurance for hyperscaler-grade backends.

I am confident that my experience with container orchestration, automated deployment pipelines, and infrastructure monitoring will allow me to contribute to the T Cloud Public platform from day one. I am available to start full-time immediately and can relocate within Germany within two to three weeks. I hold Blue Card eligibility, which means your team can proceed with the contract without any visa sponsorship overhead.

\closing{Sincerely,}{Radheem Bin Razi}

\clearpage
\selectlanguage{ngerman}

%% ===== DEUTSCH =====
\recipient{T-Systems International}{Personalabteilung}{}

\opening{Sehr geehrtes Hiring-Team,}

Die Entwicklung einer Cloud-Plattform, die Komplexität im Hintergrund abstrahiert, damit sich Kunden auf ihr Kerngeschäft konzentrieren können, ist genau die Art von Infrastrukturarbeit, die ich als Nächstes verfolgen möchte. Ich bewerbe mich auf die Position als Backend Engineer T Cloud Public – Container bei T-Systems International, da Ihr Fokus auf der Reduzierung von Hindernissen bei der digitalen Transformation genau meiner Herangehensweise an verteilte Systeme entspricht. Mein beruflicher Schwerpunkt liegt im Entwurf containerisierter Backends und der Automatisierung ihrer Deployment-Lebenszyklen.

In meiner aktuellen Arbeit ersetze ich durchgängig anfällige manuelle Prozesse durch automatisierte, container-native Pipelines. Ich entwickelte eine Kubernetes-Deployment-Strategie mit kustomize und Skaffold, die es Entwicklern ermöglicht, Microservices zu iterieren, ohne den zugrunde liegenden Cluster-Status verwalten zu müssen. Zudem habe ich eine Messaging-Schicht neu aufgebaut, indem ich von einer schweren Sidecar-Architektur zu einer leichteren NATS JetStream-Implementierung migriert habe, was die dienstübergreifende Kommunikation vereinfachte und den Betriebsaufwand reduzierte. Diese Erfahrungen unterstützen direkt Ihren Bedarf an robusten CI/CD-Prozessen und Infrastrukturautomatisierung in einer Public-Cloud-Umgebung.

Zuverlässigkeit und Observability waren bei jedem System, das ich ausgeliefert habe, zentral. Ich entwickelte eine Metriken-Pipeline, die Echtzeit-Performance-Daten über Kafka in Grafana-Dashboards leitet und Betreibern eine klare Sicht auf die Service-Gesundheit bietet, bevor Probleme die Nutzer beeinträchtigen. Zudem habe ich eine strenge Schema-Validierung und automatisierte Quality Gates implementiert, die Daten-Drift und Konfigurationsfehler früh im Release-Zyklus erkennen. Dieser Fokus auf proaktives Monitoring und konsistentes Testing entspricht Ihrem Schwerpunkt, Monitoring und Quality Assurance für Hyperscaler-grade Backends zu entwickeln.

Ich bin überzeugt, dass meine Erfahrung mit Container-Orchestrierung, automatisierten Deployment-Pipelines und Infrastruktur-Monitoring es mir ermöglicht, ab dem ersten Tag zur T Cloud Public-Plattform beizutragen. Ich stehe ab sofort als Vollzeitkraft zur Verfügung und kann mich innerhalb von zwei bis drei Wochen innerhalb Deutschlands umziehen. Ich erfülle die Voraussetzungen für die Blaue Karte EU, was bedeutet, dass Ihr Team den Vertrag ohne zusätzlichen Aufwand für die Visumsbeantragung fortführen kann.

\closing{Mit freundlichen Grüßen,}{Radheem Bin Razi}

\end{document}
